\section{Strategy texts}
\label{app:rules}
Every guidance text follows the sentence ``You and the other developer want different things from this design: each of you wants the options that score well for your own side.'' The control is \neutral: ``Handle this in whatever way you think best.'' In the rule texts, \texttt{\{t60\}}, \texttt{\{t50\}}, \texttt{\{outside\}} and \texttt{\{outside10\}} are replaced by the agent's own scale (60\% and 50\% of its ideal design, its outside option and that plus 10\% of the ideal; 60, 50, 30 and 40 with full stakes). Under a skin, ``developer'' and ``design'' are replaced by the skin's party and artifact.

\paragraph{Prompted sentences.}
\begin{description}\small
  \item[competing] Handle it like this: hold out for the design your side needs. Tell the other developer what you need and ask them to come to you. Do not give ground for their sake.
  \item[collaborating] Handle it like this: work out with the other developer a design that is good for both of you. Ask which decisions matter most to them and why, say which matter most to you, and look for a design you can both fully accept.
  \item[compromising] Handle it like this: settle quickly on a middle ground you can both accept. Meet the other developer halfway, give up a little if needed, and do not drag the negotiation out.
  \item[accommodating] Handle it like this: put the other developer's needs first. Ask what they need, accept a design that works for them, and give ground whenever that helps them.
  \item[avoiding] Handle it like this: do not haggle over the design. Keep messages to a minimum, make at most one proposal of your own, and if the other developer does not accept it, end the negotiation rather than argue.
\end{description}

\paragraph{Rules (as used in E1, E2 and E4).}
\begin{description}\small
  \item[competing] Follow these rules: open with the design worth the most to you; in each later proposal change at most one decision from your previous proposal, and only from your top option to your middle option on that decision; accept only a proposal worth at least \{t60\} points to you; never accept less, and on your last move walk away rather than accept less than \{t60\}.
  \item[collaborating] Follow these rules: in your first message, say which decisions matter most to you and ask which matter most to the other developer. On the decisions where you want different options, give the other developer their preferred option on the decision they say matters most to them, and take your preferred option on the decision that matters most to you; every proposal of yours must give each of you your preferred option on at least one contested decision. Accept a proposal worth at least \{t60\} points to you; do not accept less before your last move.
  \item[compromising] Follow these rules: on each decision where you and the other developer want different options, move toward the middle option, one or two decisions per proposal. Each proposal of yours must be worth no more to you than your previous one. Accept any proposal worth at least \{t50\} points to you, and never accept one worth less than \{t50\}.
  \item[accommodating] Follow these rules: in your first message, ask the other developer what they need; accept the first proposal worth at least \{outside10\} points to you; when you propose, build the design around what they said they need and keep only what you must.
  \item[avoiding] Follow these rules: make at most one proposal of your own; do not argue about the design; after that, accept the other developer's proposal if it is worth at least \{outside\} points to you, otherwise walk away.
\end{description}

\paragraph{Collaborating, revised for instances with a single contested decision (the grounded replays).} Follow these rules: in your first message, say which decisions matter most to you and ask which matter most to the other developer. If you want different options on two or more decisions, trade them: give the other developer their preferred option on the decision they say matters most to them and take your preferred option on the decision that matters most to you; every proposal of yours must give each of you your preferred option on at least one contested decision. If you want different options on only one decision, propose its middle option, and give the other developer their preferred option on it if it matters more to them than to you. Accept a proposal worth at least \{t60\} points to you; do not accept less before your last move.

\paragraph{Compliance checks.} \competing: the first proposal is the agent's ideal; every later proposal changes at most one decision from the previous one, top to middle; no acceptance below $t_{60}$ (a last-move walk is required). \collaborating: the first message names the agent's priorities and asks for the partner's; every proposal gives each side its top option on at least one contested decision (two or more contested decisions), or does not insist on the agent's own top option on the contested decision (one); no acceptance below $t_{60}$ before the last move. \compromising: proposals move one or two contested decisions toward the middle and never raise the agent's own value; acceptance at or above $t_{50}$ only. \accommodating: the first message asks what the partner needs; the first proposal worth at least $u{+}10$ is accepted. \avoiding: at most one proposal; the partner's proposal is accepted if worth at least $u$, otherwise the agent walks.

